\begin{xltabular}{\textwidth}{| >{\bfseries}l | X |}

% --- INTESTAZIONE PRIMA PAGINA ---
\caption{Caso d'Uso: Gestione di un Alert di Qualità} \label{tab:uc_rating_alert} \\
\hline
\rowcolor{gray!15} Campo & \textbf{Contenuto} \\ \hline
\endfirsthead

% --- INTESTAZIONE PAGINE SUCCESSIVE ---
\multicolumn{2}{c}{{\tablename\ \thetable{} -- Continua dalla pagina precedente}} \\
\hline
\rowcolor{gray!15} Campo & \textbf{Contenuto} \\ \hline
\endhead

% --- PIÈ DI PAGINA (se si spezza) ---
\hline
\multicolumn{2}{r}{\textit{Continua nella pagina successiva...}} \\
\endfoot

% --- PIÈ DI PAGINA FINALE ---
\endlastfoot

% --- CONTENUTO ---
Nome & Gestione di un Alert di Qualità \\ \hline
Livello & Obiettivo Utente \\ \hline
Descrizione & L'Energy Manager esamina una segnalazione aperta automaticamente dal sistema su una colonnina la cui qualità percepita è degradata, e decide se sospenderla oppure archiviare la segnalazione. \\ \hline
Pre-condizioni & \parbox[t]{\linewidth}{%
L'Energy Manager ha effettuato l'accesso.\\
Esiste almeno un \texttt{RatingAlert} in stato \texttt{PENDING}. La segnalazione non nasce da una singola recensione negativa: è aperta dal \texttt{RatingService} quando la media della colonnina scende sotto 2.0 con almeno 5 valutazioni. Per una stessa colonnina non può esistere più di una segnalazione aperta.
} \\ \hline
Post-condizioni & La segnalazione passa a \texttt{RESOLVED\_SUSPENDED} oppure a \texttt{RESOLVED\_IGNORED}, con la nota del Manager registrata. Solo nel primo caso la colonnina passa a \texttt{SUSPENDED}; archiviando, il suo stato resta invariato. \\ \hline
Attori & Energy Manager \\ \hline
Unit Test & \texttt{RatingServiceTest} \\ \hline
Flusso Principale &
\begin{enumerate}[leftmargin=*, nosep]
    \item L'Energy Manager apre la schermata di gestione delle segnalazioni di qualità.
    \item Il sistema mostra le sole segnalazioni in stato \texttt{PENDING}, con nome della colonnina, media attuale delle valutazioni, indicatore di gravità e tempo trascorso dall'apertura.
    \item L'Energy Manager seleziona una segnalazione e sceglie una delle due azioni possibili:
    \begin{itemize}[leftmargin=*, nosep]
        \item \textbf{Archivia la segnalazione:}
        \begin{itemize}[leftmargin=*, nosep]
            \item[4a.] Inserisce una nota di motivazione, obbligatoria, e conferma.
            \item[5a.] Il sistema porta la segnalazione a \texttt{RESOLVED\_IGNORED} e ne registra la nota. Lo stato della colonnina non viene toccato.
        \end{itemize}
        \item \textbf{Sospendi la colonnina (\texttt{<<extend>>}):}
        \begin{itemize}[leftmargin=*, nosep]
            \item[4b.] Inserisce la motivazione della sospensione, obbligatoria, e conferma.
            \item[5b.] Il sistema porta la segnalazione a \texttt{RESOLVED\_SUSPENDED} e la colonnina a \texttt{SUSPENDED}, in un'unica transazione.
        \end{itemize}
    \end{itemize}
    \item La segnalazione scompare dalla coda di quelle aperte.
\end{enumerate} \\ \hline
Flussi Alternativi &
\textbf{2. Segnalazione già gestita:} \newline
2.1. Se nel frattempo la segnalazione non è più in stato \texttt{PENDING}, ad esempio perché chiusa da un altro Manager, il sistema rifiuta l'operazione. \newline
2.2. La decisione non spetta al servizio ma all'entità di dominio: \texttt{RatingAlert} applica la risoluzione solo se la segnalazione è ancora \texttt{PENDING}, quindi una già chiusa non può essere risolta una seconda volta. \newline\newline
\textbf{4. Annullamento:} \newline
4.1. Se il Manager chiude il dialogo o lascia vuota la nota, l'operazione viene abbandonata e nessuna modifica viene salvata. \newline\newline
\textbf{Nota sull'atomicità:} \newline
\textbullet\ Nella sospensione, l'aggiornamento della segnalazione e quello della colonnina appartengono a un'unica transazione: se una delle due scritture fallisce, il \texttt{rollback} annulla anche l'altra, evitando una colonnina sospesa con la segnalazione ancora aperta, o viceversa. \\ \hline

\end{xltabular}
